\documentclass[../main.tex]{subfiles}
\begin{document}
%==========================================TIÊU ĐỀ TẬP========================================
\begin{table}[h]
    \begin{adjustwidth}{-1cm}{}
        \begin{tabular}{>{\centering\arraybackslash}p{6cm}|>{\raggedright\arraybackslash}p{12.5cm}}
            \multirow{5}{6cm}
            % Cột bên trái: ÉPISODE và số tập
            {\rainbowcolor
                {\\[-40pt]
                    \flaregothic\fontsize{20pt}{20pt}\selectfont \centering
                    \color{eptype}JOUR\color{black}\\[12pt]
                    \hbrk\fontsize{72pt}{72pt}\selectfont
                    \color{epnum}2
                    \\[18pt]
                    % \flaregothic\fontsize{14pt}{14pt}\selectfont
                    % \color{epattr}BIENVENUE, \uline{\color{Banpire} Mel} !
                    % \flaregothic\fontsize{18pt}{18pt}\selectfont
                    % \color{epattr}ÉPISODE PERDU\\
                    \unrainbowcolor}
                \color{black}}
             &
            % Dòng thứ nhất: tên tập tiếng Nhật
            {\fotskip\fontsize{18pt}{18pt}\selectfont \ruby{万}{ばん}\ruby{策}{さく}\ruby{尽}{つ}きる
            }           \\[6pt]
            % Dòng thứ hai: tên tập tiếng Anh
             & {\cafeteria\fontsize{18pt}{18pt}\selectfont We Messed Up\dots}                        \\[4pt]
            % Dòng thứ ba: tên tập tiếng Việt
             & {\vnmsans\fontsize{18pt}{18pt}\selectfont Cuộc đi dạo phiên bản cạn túi}              \\[8pt]
            % Dòng thứ tư: Nhân vật xuất hiện
             & {\sen{\fontsize{14pt}{14pt}\selectfont Track used: Ice Sheet (from Rolling Fanmade)}}
            \\[6pt]
            % Dòng cuối cùng: Thời gian diễn ra của tập (thường sẽ là ngày phát hành)
             & {\continuum\fontsize{16pt}{16pt}\selectfont Ngày phát hành: 02/05/2021}               \\[18pt]
        \end{tabular}
    \end{adjustwidth}
\end{table}
%=========================================TRUYỆN CHÍNH========================================
\color{black}

\pov{Hikawa Suzuki}

Ngày thứ hai mà chúng tôi lạc ở thế giới Nijisanji.

Chúng tôi quyết định tranh thủ thời tiết đẹp để ra ngoài hít thở không khí trong lành và thư giãn sau ngày đầu tiên đầy sự nhốn nháo và có phần mất kiểm soát.

Bầu trời trong xanh như vừa được lau sạch, chỉ lác đác vài cụm mây trắng trôi chầm chậm. Không khí có chút se lạnh dễ chịu, vừa đủ để tôi có thể mặc áo khoác mỏng mà vẫn cảm thấy ấm áp. Hoa anh đào đang nở rộ khắp hai bên bờ sông, từng cánh hồng phớt theo gió bay nhẹ, như tuyết xuân rơi.

Tôi, Onii-sama và hai chị VTuber mà tôi vẫn chưa kịp làm quen kỹ đang đi dạo trên con đường lát đá chạy dọc bờ sông, chỉ cách công ty một đoạn ngắn. Không gian yên bình đến lạ. Tiếng chim hót líu lo hòa với tiếng gió xào xạc của cánh hoa anh đào làm tôi cứ ngỡ mình đang mơ.

Trong khi đó, Game-san thì đang nghỉ ở nhà sau một ngày dài mệt mỏi, trong khi Onee-sama thì ở lại phòng ký túc để dọn dẹp lại chỗ ở và chuẩn bị bữa trưa cho chúng tôi.

Nhưng chằng ngờ, chính buổi sáng tưởng như bình thường ấy lại mở đầu cho một chuỗi sự kiện mà không ai trong chúng tôi có thể đoán trước. Tất cả chỉ vì một buổi đi dạo tưởng như bình yên kia.

\img{nj-2a}

\chrint

\begin{wrapfigure}[29]{l}{7cm}
    \includegraphics[width=7cm]{../../../../Image/Book?/saku_model.png}
\end{wrapfigure}

Một trong hai chị đi cùng chúng tôi là Sasaki Saku (\ruby{笹}{さ}\ruby{木}{さき}\ruby{咲}{さく}, \ruby{Xa}{Thế}-\ruby{xa-ki}{Mộc}\ \ruby{Xa-cư}{Tiếu}). Chị là một thành viên thuộc nhóm GAMERS - theo như Onii-sama đề cập là cũng giống như ở bên \hll, nhưng đã bị giải thể. Chị ấy học năm hai cấp ba, có gốc Kansai, và rất điềm tĩnh.

Saku-nee có sở thích là chơi game - tôi nghĩ đó là lý do vì sao chị ấy ở trong nhóm GAMERS đó - và có một tình yêu với gấu trúc.

Chị ấy rất thích các trò chơi điện tử, đặc biệt là những game có yếu tố cạnh tranh hoặc cần sự nhanh nhạy, như Mario Kart hay các game thể thao điện tử. Tôi nghe nói chị từng tổ chức nhiều buổi stream chơi game cùng các thành viên khác, và mỗi lần như vậy không khí đều rất náo nhiệt. Ngoài ra, Saku-nee còn nổi tiếng với những màn ``rage quit'' (tức giận vì thua quá nhiều) hài hước, mỗi khi thua game là chị ấy sẽ hét lên hoặc giả vờ giận dỗi, nhưng thực ra lại rất dễ thương.

Theo như những gì mà mọi người ở Nijisanji nói, chị ấy là người dễ nổi nóng và thường xuyên cáu gắt ở những trò mà chị ấy chơi. Tôi cũng được các anh chị ở đó nói là chị ấy từng rời công ty vì bất hòa trong vấn đề về game, nhưng rồi sau khi mọi chuyện được giải quyết, chị ấy đã quay trở lại.

Tôi còn nghe kể rằng Saku-nee rất yêu động vật, đặc biệt là gấu trúc. Chị ấy thường xuyên nhắc đến gấu trúc trong các buổi trò chuyện, thậm chí còn sưu tầm rất nhiều đồ vật liên quan đến gấu trúc nữa. Có lần tôi thấy chị ấy ôm một con gấu trúc bông rất to, trông vừa buồn cười vừa đáng yêu.

Theo hiểu biết sơ khai của Onii-sama, Saku-nee cũng là người hay chơi chung với Aqua-nee và Okayu-nee ở bên \hll\ từ báo cáo của hai chị ấy. Chị ấy và Yuika-nee cũng là bạn chí cốt của nhau nữa, tới mức hầu như không bao giờ tách nhau.

Tuy là một người dịu dàng, nhưng đôi lúc Saku-nee cũng xảy ra mâu thuẫn với những đồng nghiệp nam, nhất là từ Banderas-san\footnote{Sẽ xuất hiện ở {\flaregothic JOUR \hbrk 5}}, dù tôi biết ông ấy luôn cởi mở với mọi người, kể cả chúng tôi.

Saku-nee có mái tóc hồng, ngắn tới gần vai, có một chỏm tóc ngố (mà bình thường người ta không thấy do chị ấy hay đội mũ trùm) và có đôi mắt hồng. Chị còn có một cái răng nanh lộ ra ở mép trái miệng. Chị mặc bộ trang phục giống như một nữ sinh trung học, với áo sơ-mi trắng có cổ bị áng hồng và được cài cúc chỉn chu nhưng không sơ-vin, đeo một chiếc nơ đỏ trước ngực, váy ngắn xếp ly màu xanh lá họa tiết ca-rô, đi giày thể thao xanh sẫm cao tới che mắt cá chân và mặc áo khoác trùm đầu hình gấu trúc.

\begin{wrapfigure}[29]{r}{7.5cm}
    \includegraphics[width=7.5cm]{../../../../Image/Book?/yuika_model.png}
\end{wrapfigure}

Đi bên cạnh Saku-nee là Shiina Yuika (\ruby{椎}{しい}\ruby{名}{な}\ruby{唯}{ゆい}\ruby{華}{か}, \ruby{Si-i}{Chuy}-\ruby{na}{Danh}\ \ruby{Giui}{Duy}-\ruby{ca}{Hoa}), như đã nói ở trên, là bạn thân của chị ấy, và cũng từng chung nhóm GAMERS luôn, cho tới khi nhóm bị giải thể. Tôi không dám chắc chị ấy bao nhiêu tuổi, nhưng tôi nghĩ chị ấy đâu đó ngang tuổi với cô bạn thân của mình.

Tôi nghe nói Yuika-nee rất nổi tiếng với khả năng kể chuyện ma và những câu chuyện kỳ bí. Chị ấy thường tổ chức các buổi stream kể chuyện rùng rợn, khiến không khí trở nên hồi hộp và cuốn hút. Mỗi lần như vậy, mọi người đều chăm chú lắng nghe, có người còn bảo rằng chị ấy có thể khiến cả những người gan dạ nhất cũng phải rùng mình. Tôi không có nhiều trải nghiệm về chuyện đó, nên tôi cũng không biết chắc.

Ngoài ra, Yuika-nee còn rất thích các trò chơi điện tử, đặc biệt là những game có yếu tố phiêu lưu hoặc giải đố. Chị ấy thường xuyên hợp tác với Saku-nee trong các buổi stream chơi game, tạo nên một cặp đôi ăn ý và hài hước. Tôi nghe kể rằng hai chị ấy từng có một sê-ri chơi game cùng nhau rất nổi tiếng, thu hút rất nhiều người xem và bình luận sôi nổi.

Tính cách của Yuika-nee khá đặc biệt. Chị ấy có vẻ ngoài điềm tĩnh, nhưng lại rất dễ bị cuốn vào những điều kỳ lạ hoặc những trò đùa của bạn bè. Mọi người ở Nijisanji thường trêu rằng chị ấy là ``người dễ bị ma nhập nhất công ty'', vì mỗi lần có chuyện gì kỳ lạ xảy ra, Yuika-nee luôn là người đầu tiên bị nghi ngờ. Tuy vậy, chị ấy luôn đón nhận mọi chuyện với thái độ vui vẻ và không hề tỏ ra sợ hãi.

Tôi cũng được biết Yuika-nee rất quan tâm đến bạn bè và đồng nghiệp. Dù đôi khi có vẻ lơ đãng hoặc tự lập, nhưng chị ấy luôn sẵn sàng giúp đỡ mọi người khi cần thiết. Sự dịu dàng và hài hước của Yuika-nee khiến chị ấy trở thành một người bạn đáng tin cậy và được mọi người yêu mến trong tập thể Nijisanji.

Yuika-nee có đôi mắt đỏ và mái tóc trắng ánh hồng được buộc đuôi ngựa thấp ở phía sau bằng một cái ruy-băng đỏ. Bộ đồng phục của chị ấy dường như khá tương đồng với Saku-nee, cũng là áo sơ-mi trắng có cổ và nơ đỏ trước ngực, chỉ khác là chị ấy mặc áo len bên ngoài màu be và váy xếp ly màu đỏ họa tiết ca-rô. Chị ấy đi tất trắng cao đến đầu gối và đi giày lười.

\chrint

``Thật là một ngày hoàn hảo để đi ra ngoài,'' anh tôi gác tay sau gáy, ngước nhìn lên trời và thở dài một hơi nhẹ nhõm. ``Chỗ này cũng khá là vắng vẻ nữa. Thật thư giãn.''

``E-Em cũng thấy vậy, Onii-sama,'' tôi nắm tay anh lại, đung đưa qua những tán cây hoa anh đào đang rơi.

``Quả thật đây đúng là tiết trời mùa xuân rồi,'' chị đeo áo khoác có mũ gấu trúc nói.

``Cơ mà\dots'' cô chị mặc đồng phục học sinh trung học lên tiếng, cảm thấy có gì đó bất thường. ``Tại sao chúng ta lại không thể cử động được nhỉ?''

Thật ra\dots\ chỉ có tôi và Onii-sama là đang thực sự đi bộ. Hai chị còn lại, Yuika-nee-nee và chị Saku-nee thì\dots\ làm sao nhỉ, có vẻ như không thể cử động được.

Dù họ có thể trò chuyện với chúng tôi (thật ra là có người lồng tiếng cho họ từ phòng thu âm ở đâu đó), nhưng thân thể thì cứ như mấy tấm biển quảng cáo biết nói vậy. Họ không thể chớp mắt, không thể quay đầu, thậm chí đứng cũng chỉ có đúng một kiểu dáng. Mỗi khi cảnh thay đổi, thì hình họ cũng\dots\ thay theo luôn. Cứ như bị ``cắt'' rồi ``dán'' vào vậy.

\begin{figure}[H]
    \centering
    \includegraphics[width=12.5cm]{../../../../Image/Book?/nj-2b.png}
    \caption{Đúng, bạn không nhìn nhầm đâu. Họ đã giữ nguyên dáng đứng như thế từ ảnh trước.}
\end{figure}

Tôi vẫn chưa hiểu lắm, nên khẽ kéo tay áo Onii-sama và hỏi nhỏ:

``\hl{Onii-sama\dots{} tại sao hai chị ấy lại không thể cử động được ạ?}''

Anh ấy mỉm cười, nhìn thẳng ra dòng sông đang lấp lánh nắng chiều. ``Anh cũng thắc mắc giống em thôi. Nhưng chắc là em chưa nghe vụ cạn túi hả?''

``Cạn\dots\ túi?'' Tôi nghiêng đầu, không hiểu. Anh chưa kịp trả lời thì Saku-nee đã lên tiếng, dù vẫn đứng bất động như tượng gỗ: ``Có vẻ như ngân sách cho mùa này đã bị dùng hết vào tập đầu tiên mất rồi.''

``Ể\dots'' tôi tròn mắt. Anh tôi đổ mồ hôi hột. ``Anh tưởng cái đó chỉ nằm trong đầu Toya thôi chứ? Một kịch bản tưởng tượng thôi mà?''

Saku-nee vẫn giữ nguyên tư thế thả lỏng ấy, nhẹ nhàng chen vào: ``Ừ, cái đó cũng được tính vào chi phí sản xuất rồi.''

Yuika-nee sau đó đột ngột lên tiếng sau vấn đề, giọng như thể vừa nghe thấy điều vô lý nhất trong ngày: ``Bà đang nói cái quái gì thế?''

Khung hình chuyển qua khuôn mặt chị ấy với đôi mắt trợn tròn (dĩ nhiên là trợn sẵn, vì model của chị không thể cử động thật sự). Tôi suýt bật cười, nhưng vẫn cố giữ cho nét mặt nghiêm túc vì lịch sự. Saku-nee thì vẫn giữ nguyên tư thế dang tay từ nãy, khẽ thở dài: ``Dù sao thì\dots\ chúng ta có mấy dịp được cử động như hôm đó đâu.''

Người bạn của chị ấy quay sang tạo dáng huơ tay, dù thật ra đó chỉ là hiệu ứng ghép hình chuyển động đơn giản. ``Hy vọng là lần tới họ dùng ngân sách hợp lý hơn\dots''

\begin{figure}[H]
    \centering
    \begin{subfigure}[t]{0.45\textwidth}
        \centering
        \includegraphics[width=8cm]{../../../../Image/Book?/nj-2c1.png}
    \end{subfigure}
    \quad
    \begin{subfigure}[t]{0.45\textwidth}
        \centering
        \includegraphics[width=8cm]{../../../../Image/Book?/nj-2c2.png}
    \end{subfigure}
    \quad
    \begin{subfigure}[t]{0.45\textwidth}
        \centering
        \includegraphics[width=8cm]{../../../../Image/Book?/nj-2c3.png}
    \end{subfigure}
\end{figure}

Tôi khẽ nghiêng đầu, rồi dè dặt hỏi: ``Ư-ưm\dots\ Tức là mọi người không biết trước khi quay hình là sẽ được cử động bao nhiêu lần ạ?''

``Không ai biết hết, Suzu-chan,'' Yuika-nee đáp, giọng dịu hơn. ``Cứ như chơi xổ số vậy.''

Tôi quay sang nhìn hai chị, rồi nhìn Onii-sama. Càng lúc tôi càng rối. ``Vậy là\dots\ mấy cảnh ngày hôm qua, lúc nhảy múa, hát hò\dots\ thật sự là có thật luôn ạ?''

``Anh cũng không chắc nữa,'' Onii-sama nói, mặt vẫn chưa hết ngơ ngác. ``Bên này mọi thứ đều hoạt động theo kiểu bất chấp logic, chả khác bên \hll\ là bao.''

Tôi thở ra một hơi, tay ôm lấy túi bánh mì tôi đang cầm. Tự dưng thấy tội cho hai chị ấy ghê.

``Anh đoán rằng là do hai em ấy dùng model quá đà hay là do quản lý tài chính bên Nijisanji có vấn đề,'' anh tôi lên giả thuyết. ``Bên anh dùng model hầu như hàng tuần, mà vẫn chưa bị đơ cứng lần nào. Ít nhất bọn anh vẫn còn bài Live2D dùng sơ cua nữa.''

Tôi chẳng hiểu gì cả, chỉ khẽ gật đầu như đã hiểu. Cả nhóm chỉ biết thở dài và tiếp tục đi bộ bên bờ sông.

Rồi, bỗng Saku-nee gọi: ``Mà này, Kuon-senpai?''

``Gì thế?'' Onii-sama quay sang, mắt vẫn đang ngắm hàng hoa anh đào.

``Em nghe nói là anh đang làm ở \hll\ đúng không? Thế thì\dots''

``\dots Bọn anh thường quản lý chi tiêu như thế nào vậy?'' Yuika-nee tiếp lời.

Onii-sama đưa hai tay ra sau gáy, vừa đi vừa nghĩ. ``Hmm\dots\ Bình thường thì anh hoặc A-san - chị Yuujin A ấy, người anh làm cùng - sẽ giải quyết mấy khâu kiểu đó.''

Tôi liền chen vào, hơi nâng giọng vì thấy hơi giống câu của chú Game ngày hôm qua: ``Vậy các anh chắc phải chi tiêu tiết kiệm lắm nhỉ?''

Anh lắc đầu: ``Thật ra thì họ chi tiêu khá hào phóng. Nhưng được cái là có chừng mực, và luôn biết rõ ngân sách sẽ được dùng cho những gì.''

``Thế ạ?'' cả hai chị đồng thanh, ánh mắt (cố định) như sáng lên.

``Ừ. Cũng nhờ bọn anh hướng dẫn cụ thể. Có hệ thống cả.''

``Ra thế\dots'' hai người họ ``ưm'' một tiếng, vẻ mặt vẫn không thay đổi, nhưng tông giọng cho thấy họ đang thật sự suy nghĩ về điều đó.

Tôi ngước nhìn anh, tò mò hỏi tiếp: ``Thế ngân sách bên anh thường dùng cho những việc gì ạ?''

Onii-sama thở dài một tiếng, kiểu thở dài đặc trưng mỗi khi nhắc đến chỗ làm. ``Trước đây thì phải chi nhiều cho sửa chữa văn phòng.''

``Sửa chữa?'' tôi hỏi lại.

``Ừ. Cửa kính, tường, trần nhà\dots\ từng có đợt mấy đứa đập vỡ cửa sổ tầng sáu bằng một quả bóng yoga. Có lần khác, gãy cả sàn nhà vì nhảy múa quá sung. Nhiều lúc nhiều cửa sổ bị vỡ quá thường xuyên tới mức một tuần thay hai, ba lần, khiến cho A-san phát điên.''

Tôi nuốt nước bọt. ``Đó là công ty VTuber đúng không ạ?''

``Ừ. Nhưng mà yên tâm. Từ khi anh lắp hệ thống tái tạo vật chất ba tháng trước, thì gần như không cần tốn khoản đó nữa.''

``Thế\dots\ ngân sách được chuyển qua mục nào khác ạ?'' tôi hỏi, vì thật sự tò mò.

``Mục 3D,'' anh đáp không cần suy nghĩ. ``Model, animation, render - tất cả đều ưu tiên. Nhờ đó mà anh, Kaigou-chan với em mới được cử động mượt mà thế này đó.''

Tôi nhìn xuống đôi tay mình, rồi nhấc một chân lên thử nhún nhảy, đúng là cảm giác có thật. Có cảm giác như tôi được là chính mình. Anh tiếp tục: ``Dù em không thuộc \hll, nhưng vì đang đi cùng anh, nên được mượn hiệu ứng. Cũng giống như Ryuuhou, hay bố mẹ anh, hay bác Hijaki-san thôi.''

``Em có lẽ là được đi ké rồi\dots'' tôi lẩm bẩm, rồi chợt bật cười.

Đúng lúc ấy, Saku-nee lên tiếng, như muốn đổi chủ đề: ``Mà này, chúng ta đang đi dạo đấy. Sao không tạm gác chuyện ngân sách sang một bên đi?''

``Đồng ý!'' Yuika-nee giơ tay (dù tay chị vẫn không nhúc nhích).

Onii-sama gật đầu: ``Ừ, có lý. Mùa xuân chỉ có một lần mỗi năm.''

Tôi nhìn theo cánh hoa anh đào rơi xuống trước mặt, tay vô thức đưa ra đón lấy.

Một giây đó\dots\ tôi cảm thấy mình thật sự đang sống.

\ct

Những hàng anh đào ven đường lúc này đã bung nở mạnh nhất, từng chùm hoa xòe ra rực rỡ như thể muốn khoe hết vẻ đẹp cuối cùng trước khi mùa xuân trôi qua. Thi thoảng tôi lại thấy vài nhánh lá non màu xanh lục bắt đầu chen lên giữa sắc hồng dịu. Thời tiết thật tuyệt vời: trời mát, nắng nhẹ, không quá chói. Một ngày hoàn hảo để ra ngoài, mà không cần lý do gì cả.

\img{nj-2d}

Chúng tôi cứ thế tiếp tục đi bộ, nhịp chân đều đều trên lối đi rải sỏi ven sông. Xung quanh yên tĩnh, nhưng không đến mức buồn tẻ. Cứ như cái kiểu yên lặng mà mình vẫn nghe thấy chim hót, lá rơi và tiếng gió, nhưng lại không thấy ai làm phiền gì cả.

``Ngẫm lại thì\dots'' Yuika-nee lên tiếng, mắt nhìn lên những tán cây rung rinh trong gió. ``Thời tiết hôm nay đẹp thật nhỉ?''

Onii-sama cười khì. ``Bọn mình vừa mới nói vậy mà. Nhưng\dots\ đúng là trời đẹp thật.''

Tôi nhìn họ, rồi không kìm được mà buột miệng: ``Ư-ừm\dots\ Nhưng mà trời đẹp như vậy mà hai chị vẫn không được cử động, thì cũng hơi buồn nhỉ.''

Câu nói khiến cả nhóm khựng lại một chút.

``Ừ th\dots\'' Saku-nee thở ra một tiếng. ``Cơ thể chúng tôi đang ở trong studio mà. Model không được cấp phép di chuyển thì đành đứng yên thôi.''

``Em nghĩ là có thể do chị dùng model quá đà hôm qua đó ạ,'' tôi vừa nói vừa rụt cổ vì sợ vô lễ. ``À, em xin lỗi nếu nói vậy là không phải\dots''

Chị ấy chỉ cười khổ. ``Cái đó\dots\ có thể lắm. Chị ghen tỵ vì em có thể di chuyển thoải mái đấy.''

``Hoa anh đào nở rộ ghê\dots'' Yuika-nee nói, như để chuyển đề tài, mắt vẫn không rời khỏi hàng cây hai bên bờ sông.

``V-Vâng\dots\ Giá mà năm sau mọi người vẫn còn ở đây để ngắm hoa nhỉ,'' tôi lẩm bẩm, không biết mình đang nói cho ai nghe.

``Phải ha,'' Saku-nee tiếp lời. ``Chán thật đấy. Mấy bông đẹp nhất cũng sắp tàn rồi\dots''

\begin{figure}[H]
    \centering
    \begin{subfigure}[t]{0.45\textwidth}
        \centering
        \includegraphics[width=8cm]{../../../../Image/Book?/nj-2f1.png}
    \end{subfigure}
    \quad
    \begin{subfigure}[t]{0.45\textwidth}
        \centering
        \includegraphics[width=8cm]{../../../../Image/Book?/nj-2f2.png}
    \end{subfigure}
\end{figure}

``Thôi nào,'' Onii-sama khẽ vỗ vào vai hai chị ấy, nở một nụ cười rất thân thiết: ``Không năm nay thì năm sau đi cũng được mà.''

Tôi cúi đầu, cười nhẹ. Rồi nhỏ giọng nói thêm: ``Nhưng nếu năm sau vẫn không đủ ngân sách cho model thì\dots\ chắc lại đứng im ngắm từ xa thôi ạ\dots''

Cả ba người đều phá lên cười nhẹ.

\ct

Một lúc sau, Saku-nee đột nhiên nhìn lên cây, ánh mắt sáng lên như vừa nhớ ra điều gì đó.

``Mà, hoa anh đào à\dots'' chị ấy nói, rồi nghiêng đầu một chút như đang trầm tư. ``Hai người đã bao giờ nghe câu `Bên dưới gốc cây anh đào có chôn xác chết'\footnote{Ý truyện ngắn 「櫻の樹の下には」 (Bên dưới gốc anh đào) của Kaijii Motojiro (1901-1932). Về nội dung và cả tác phẩm, có thể xem tại \href{https://spiderum.com/bai-dang/Mot-truyen-kinh-di-Ben-duoi-goc-anh-dao-82a}{đây}. Câu nói của Saku chính là câu đầu tiên của tác phẩm ấy.} chưa?''

\img{nj-2g}

Cả tôi và Onii-sama cùng ngẩng đầu nhìn chị. ``Ý em là câu mở đầu từ truyện ngắn nổi tiếng ấy à?'' anh tôi hỏi. ``Có, nhưng sao?''

``Quả đúng là một cách mở đầu giật gân cho một truyện ngắn nhỉ?'' Saku-nee cười nhạt.

Tôi chưa kịp phản ứng gì thì Yuika-nee hỏi luôn: ``Tôi chưa nghe bao giờ. Là truyện kiểu huyền bí hay kinh dị vậy?''

``Không đâu,'' chị Gấu trúc lắc đầu. ``Không phải kiểu ma quái gì cả.''

Onii-sama chen vào: ``Nó thiên về triết lý hơn. Kiểu\dots\ suy ngẫm về sự sống và cái chết ấy.''

``Vậy thì nó kể về cái gì?'' Yuika-nee hỏi khi đụng trúng chủ đề tới tâm linh.

\img{nj-2h}

``Tóm gọn thì\dots'' Saku-nee ngẫm nghĩ một lúc, rồi giải thích: ``Chuyện kể về một nhân vật chính đang cố gắng hiểu tại sao hoa anh đào lại đẹp đến vậy. Và họ tin rằng, vẻ đẹp đó là nhờ\dots\ chất dinh dưỡng từ những cái xác chôn bên dưới gốc cây.''

``\dots Cái gì cơ?'' người bạn của chị bật cười, tỏ vẻ không tin lắm.

Onii-sama cười, gãi đầu: ``Anh lúc đầu đọc cũng thấy khó hiểu như em vậy.''

``Tôi biết là khi tóm tắt như vậy thì nghe nó chuối thật đấy,'' chị Saku thở dài. ``Nhưng thật ra nó giống một bài thơ hơn là truyện ngắn.''

``Ừ,'' Onii-sama gật đầu. ``Tác giả viết nó với sự thành kính và chiêm nghiệm.''

``Ư-ưm\dots\ nghe hơi mông lung\dots'' tôi nói khẽ. ``Nhưng chắc là do em chưa quen đọc mấy chuyện nặng như thế.''

``Nó cũng ngắn mà,'' chị Gấu trúc cười nhẹ. ``Chị nghĩ em đọc hết chỉ trong năm phút thôi.''

``Thật á?'' Yuika-nee ngạc nhiên.

``Không tin thì thôi,'' Onii-sama cười. ``Truyện chỉ dài hai ba trang giấy thôi.''

``Nếu em thích thì có thể thử,'' chị Saku quay sang tôi. ``Hoặc là nếu muốn, chị kể cho em nghe cũng được.''

Tôi gật đầu. ``Ư-ưm\dots\ Chắc em sẽ đọc thử sau. Em thấy\dots câu chuyện đó giống như là mình phải vừa đọc, vừa cảm nhận bằng tim vậy.''

Onii-sama nhìn tôi một lúc, rồi khẽ mỉm cười. ``Em nói đúng lắm, Suzuki.''

\ct

Đang đi dạo thong dong giữa tiết trời đẹp như mơ, thì Yuika-nee bất chợt lên tiếng, giọng đều đều như đang thẩm vấn một tội phạm văn học: ``Này, bà luôn có sở thích quái dị đến thế à? Bà có vẻ hiểu rõ mấy cái truyện rợn rợn đó phết đấy.''

Chị Saku bật cười, tay khoanh trước ngực: ``Đâu có? Tôi chỉ đọc mấy dòng trong kịch bản thôi ấy chứ.''

Onii-sama thở dài, đưa tay xoa trán: ``\dots Anh lại quên mất là mấy em vẫn đang ở trong phòng thu âm.''

Tôi cũng cười trừ, nói nhỏ: ``Chẳng trách tại sao hai chị lại rành rọt thế. Em cảm giác như hai chị đang phá vỡ bức tường thứ tư luôn ấy.''

Cả ba chị em tôi lại phá lên cười. Nhưng sau tiếng cười đó, tôi nghe Yuika-nee thở dài, lần này mang chút gì đó chán nản thật sự. ``Chúng ta thật sự có thể làm một tập anime ra hồn về mấy chuyện này không nhỉ?''

Saku-nee cũng nhìn sang Onii-sama và tôi, đang đi bộ rất tự nhiên như những nhân vật có chuyển động hoàn chỉnh, rồi lẩm bẩm: ``Nhất là khi bọn mình phải đứng bất động một chỗ, trong khi hai anh em nhà Kuon-senpai thì cứ đi tung tăng\dots''

``Tôi phản đối!'' Yuika-nee phụng phịu, hai má phồng lên. ``Cho bọn này hưởng ké với!''

Tôi bật cười, che miệng: ``Nếu mà được chia sẻ mô hình 3D thì em sẵn lòng đó.''

Onii-sama cười hiền: ``Bên anh có hệ thống máy tái tạo, nhưng cũng phải mất kha khá chi phí mới hoạt động được. Với cả đó là tài sản công ty, không phải lúc nào cũng được dùng đâu.''

Tôi thì thầm: ``Thật ra anh ấy rất hay ôm việc đấy, chứ có ai quản lý đâu\dots''

Chúng tôi còn đang nói chuyện thì bỗng Onii-sama nhíu mày, nhìn về phía trước. ``Hử? Hmm\dots''

Chúng tôi cùng quay lại nhìn theo hướng đó và đều khựng lại.

Dưới một gốc anh đào ở gần bờ sông, có một bóng người đang cặm cụi\dots\ đào bới?

``Chờ chút\dots'' chị Siêu nhiên nheo mắt. ``Có phải ai đó đang đào hố dưới gốc cây kia không?''

``Ừ,'' Onii-sama gật đầu. ``Anh cũng thấy thế.''

``Ể? Thật luôn á?'' chị Gấu trúc kêu lên. ``Khoan\dots\ có phải đó là\dots''

Cả ba chúng tôi bước lại gần hơn, mắt mở to đầy kinh ngạc khi thấy rõ gương mặt của một cô bé succubus quen thuộc đang lấm lem vì đất cát.

\chrint

\begin{wrapfigure}[16]{l}{9cm}
    \includegraphics[width=9cm]{../../../../Image/Book?/ririmu_model.png}
\end{wrapfigure}

Cô bé lém lỉnh đang tất bật đào đất ấy là Makaino Ririmu (\ruby{魔}{ま}\ruby{界}{かい}\ruby{ノ}{の}りりむ, \ruby{Ma}{Ma}-\ruby{Giới}{ca-i}-nô Ri-ri-mư), một cô bé succubus trẻ tuổi tới từ vùng đất Quỷ giới, và cũng từng trong nhóm GAMERS với Saku-nee và Yuika-nee. Tên đầy đủ của em ấy là Ririmu KissMe LovelyHeart = Lolitania (リム・キスミー・ラブリーハート＝ロリータニア) - một cái tên vừa dài vừa kỳ quặc.

Một cô bé 13 tuổi lên trần thế với nhiệm vụ gì đó có vẻ hiểm hóc và ác độc, nhưng rồi em ấy\dots\ quên mất nó, và dần dấn thân vào trò chơi điện tử. Theo như những gì Saku-nee và Yuika-nee nói, bố của em ấy là quỷ Satan, vị vua của Quỷ giới, nên không ai dám ho he. Trừ Onii-sama.

Tôi nghe nói Ririmu-chan rất nổi tiếng ở Nijisanji vì tính cách tinh nghịch và có phần quỷ quái đúng như xuất thân của mình. Em ấy thường xuyên bày ra đủ trò trêu chọc các thành viên khác, từ những trò đùa vô hại cho đến những pha ``troll'' khiến mọi người phải dở khóc dở cười. Dù nhỏ tuổi nhất nhóm, nhưng Ririmu-chan lại rất tự tin và không ngại cãi lại cả những đàn anh, đàn chị lớn tuổi hơn mình.

Ngoài ra, Ririmu-chan còn có biệt tài ``giả vờ ngây thơ''. Mỗi khi bị bắt quả tang làm điều gì đó nghịch ngợm, em ấy sẽ lập tức đổi giọng, làm mặt đáng yêu hoặc giả vờ khóc để được tha thứ. Nhưng chỉ cần mọi người lơ là một chút là em ấy lại quay về bản chất ``tiểu quỷ'' ngay lập tức. Tôi nghe kể rằng, nhiều lần các anh chị trong công ty phải ``bó tay'' trước sự lém lỉnh của em ấy.

Tôi cũng được biết Ririmu-chan rất mê các trò chơi điện tử, đặc biệt là những game có yếu tố hành động hoặc phiêu lưu. Em ấy thường xuyên hợp tác với các thành viên khác trong các buổi stream, tạo nên không khí vui nhộn và đầy tiếng cười. Dù đôi khi chơi không giỏi lắm, nhưng Ririmu-chan luôn biết cách khiến mọi người chú ý bằng những màn ``tấu hài'' rất riêng của mình.

Điều khiến tôi ấn tượng nhất là dù hay nghịch ngợm, Ririmu-chan lại rất tình cảm với bạn bè. Em ấy luôn quan tâm, động viên mọi người mỗi khi ai đó gặp chuyện buồn. Có lẽ chính sự kết hợp giữa nét tinh nghịch và tấm lòng ấm áp ấy đã khiến Ririmu-chan trở thành một thành viên đặc biệt được yêu mến trong Nijisanji.

Em ấy có mái tóc xám được cột thành hai bên như tôi, nhưng được cao gần với đỉnh đầu thay vì để thấp. Em ấy cũng có chỏm tóc ngố đỏ, có đôi mắt đỏ, tai nhọn vì là quỷ, có cặp sừng nhìn như sừng cừu và đôi cánh đỏ thẫm. Trang phục mà Ririmu-chan mặc được lấy từ bộ Games Day - theo lời của em ấy - gồm một bộ váy hai dây màu đen dài qua đầu gối một chút, có cánh tay áo tách rời dài đến cổ tay, đeo một chiếc túi xách đỏ chéo vai. Em ấy cũng đi tất hai sọc đen-trắng xen kẽ nhau dài tới đùi và đi bốt cao cổ màu đen.

\chrint

``Ririmu-chan!?''

``Chính là em ấy luôn!'' tôi xác nhận, hơi rướn người về phía trước. ``Nhưng\dots\ sao em ấy có thể di chuyển được nhỉ?''

\gif{nj-2i}{1}{24}

Yuika-nee tròn mắt: ``Có khi nào em ấy hối lộ ai đó không?''

``Ừm\dots\ cũng có thể lắm,'' Saku-nee nghiêng đầu suy đoán.

Onii-sama chỉ nhún vai. ``Anh không muốn nói xấu em ấy đâu\dots\ nhưng cái mồm nhanh nhảu của Ririmu-chan có thể là nguyên nhân đấy.''

Tôi nghĩ ngợi, rồi thì thào: ``Cũng có thể em ấy có một phương pháp nào đó mà tụi mình không biết\dots''

Dù thế nào đi nữa, cả nhóm đều nhất trí đứng lại quan sát. Trên đầu mỗi người như thể hiện rõ ba dấu hỏi chấm lơ lửng vô hình, nhưng dễ dàng cảm nhận.

``Nhưng mà\dots'' tôi lẩm bẩm. ``Rốt cuộc em ấy đang đào cái gì vậy?''

``Đi tìm kho báu?'' Onii-sama nghiêng đầu. ``Anh nghĩ vậy\dots''

``Cũng có lý,'' chị Siêu nhiên gật gù. ``Hay là ta lại gần xem thử đi.''

``Ừm. Miễn là không cản trở công việc\dots đào đất thần bí của em ấy,'' tôi đáp, vừa bước về phía Ririmu-chan, vừa thầm nghĩ: ``\textit{Không biết rồi mình sẽ phát hiện ra thứ gì, hay ai đó bị chôn thật dưới gốc cây nữa không đây\dots?}''

\ct

``Ririmu-chan ơi, em đang làm gì thế?'' chị Gấu trúc lên tiếng trước, đôi mắt tròn xoe vì tò mò. Onii-sama cũng bước lại gần, ngó xuống cái hố sâu hoắm mà cô bé đang đào rồi hỏi tiếp: ``Trông em có vẻ hớn hở lắm nhỉ?''

\img{nj-2j}

Nghe thấy tiếng gọi, cô bé ngẩng phắt đầu lên khỏi miệng hố, đôi má dính đầy đất cát, mắt sáng rực như đèn LED. ``Ah! Là Sasa và Shiishii nè! Với cả Ku-Kuon và Suzu nữa!?''

Cô bé cười toe toét rồi chỉ tay xuống cái hố như thể đang dẫn bọn tôi đi khai quật di tích khảo cổ: ``Có ai đó bảo với em là ở dưới này có kho báu đấy!''

\img{nj-2k}

Tôi liếc nhanh sang Yuika-nee và Onii-sama, và như có sự đồng điệu bất ngờ, cả bốn người chúng tôi đồng thanh lên tiếng: ``Đúng là như chúng ta đoán rồi\dots''

Cả nhóm ai nấy đều nhìn nhau bằng ánh mắt khó xử, còn Saku-nee thì lại có vẻ hơi phấn khích quá mức. ``\textbf{Ể!? Kho báu á!?}'' chị thốt lên như vừa nghe tin sét đánh giữa trời xuân.

``Ở chỗ thế này sao\dots'' Yuika-nee thì tỏ vẻ nghi ngờ. ``Nghe có vẻ như mấy giấc mơ lúc trưa hơn là chuyện thật đấy.''

Onii-sama thì nhìn quanh, ánh mắt như thể đang phân tích toàn bộ hiện trường. ``Anh nghĩ là em bị ai đó trêu rồi đấy, Ririmu. Làm gì có kho báu thật sự đâu.''

Tôi ngó qua cánh tay phải của Onii-sama, nơi màn hình điện thoại của anh ấy thường hiện ra giao diện cảm biến, nhưng không thấy gì lạ. Tôi lắc đầu, nói nhẹ: ``Cảm biến của chị không phát hiện kim loại nào ở đây cả. Nghiêm túc đấy, có lẽ em nên về nhà\dots\ học bài thì hơn.''

``Ể\dots!?'' Ririmu-chan buông chiếc xẻng rơi xuống đất, vẻ mặt như bị tạt một gáo nước lạnh. Đến một người mà Onee-sama bảo là ``ngây thơ'' cũng thừa hiểu rằng nơi này không có kho báu đâu.

Nhưng lạ thay, trong khi tôi cứ tưởng Yuika-nee sẽ gắt lên vì đào đất vô ích, thì chị ấy lại tỏ vẻ khá\dots\ bình thản. Thậm chí còn có chút hứng thú? ``Dù sao thì bọn chị cũng không có gì để làm mà. Giúp em một tay cũng được.''

``Biết đâu lại có cái gì vui vui ở dưới đấy thì sao?'' Saku-nee tiếp lời. ``Như\dots\ một cái xác chẳng hạn!''

Tôi đứng hình, mặt hơi tái mét. ``Hai chị thực sự đang bị truyện ma ám rồi\dots''

Dẫu vậy, lời cổ vũ đó lại như thổi thêm lửa cho cô bé Succubus nhỏ tuổi. Ririmu lập tức reo lên, nhặt lại chiếc xẻng như chưa từng có sự nghi ngờ nào trước đó.

``Vậy thì còn chờ gì nữa! Đào thôi, đào thôi nào\~{}!'' cô hét to, rồi quay sang\dots\ cà khịa Onii-sama: ``Mà Ku-Kuon đúng là đồ xấu tính!''

Hai chị em ``hình tĩnh'' phía sau thì cũng không bỏ lỡ cơ hội. Chị Siêu nhiên giả bộ chống nạnh: ``Anh ấy là người ngoài mà, chấp làm gì!''

``Tự dưng thấy hơi tội cho anh quá,'' tôi thì thào, nhìn Onii-sama bị ba cô gái ấy hội đồng. Nhưng anh ấy chỉ thở dài.

Cái kiểu chịu trận không phản kháng ấy quen thuộc lắm rồi. Hệt như lúc Onii-sama bị các chị bên \hll\ quậy phá rồi đổ hết lỗi lên đầu. Tôi lo lắng hỏi nhỏ: ``\hl{Anh không thấy khó chịu chút nào sao?}''

``Cứ để họ nói thôi,''anh cười khổ, ánh mắt như đã miễn nhiễm với mọi công kích. ``Với cả\dots\ họ đâu có nhìn thấy.''

Tôi nghiêng đầu. ``Ý anh là sao?''

Không trả lời trực tiếp, Onii-sama chỉ lặng lẽ giơ tay lên, rồi chỉ vào tấm bảng nhỏ đặt cách đó vài bước chân, hơi bị cỏ che khuất. Tôi ngước nhìn.

Một tấm bảng trắng, chữ đen rõ ràng:「立入禁止・掘削厳禁」– KHU VỰC CẤM. CẤM ĐÀO BỚI.

Tôi suýt ngã ngửa. ``Không thể nào\dots''

Đúng lúc đó, Ririmu=chan lại dúi vào tay tôi và Onii-sama mỗi người một cái xẻng nữa, hồ hởi như thể sắp trúng số.

Tôi thầm nghĩ: ``\textit{Thật sự thì\dots\ mình đang đào mộ cho ai vậy?}''

Còn Onii-sama? Anh ấy chỉ lặng lẽ thở dài, mắt vẫn dán vào tấm biển cảnh báo\dots\ như thể tiên đoán trước kết cục không mấy tốt đẹp.

\ct

Ba chị em đang đào hăng say, còn tôi và Onii-sama thì chỉ biết đứng cạnh, mỗi người chống một cái xẻng xuống đất, mặt méo xệch. Ririmu-chan lại chính là người chỉ huy ``chiến dịch khai quật''. Yuika-nee và Saku-nee thì cứ thế nghe theo răm rắp.

``Đây nhé, Shii-shii lo phần này nha\~{}'' Ririmu chỉ tay như thể đang chia địa bàn.

``Rồi rồi\~{}'' chị Siêu nhiên cười toe.

``Còn Sasa thì\dots\ lo đoạn này!'' em ấy quay ngoắt sang chị Saku, không để ai phản kháng.

``OK\~{} Đào là nghề của mình mà!'' chị Gấu trúc hăng hái đến lạ.

Tôi nhìn Onii-sama, thấy anh chỉ cười trừ rồi giơ tay tỏ vẻ từ chối. ``Ờm, thôi. Anh không muốn phải nộp phạt mấy nghìn yên chỉ vì mấy trò vô bổ này đâu.''

``Xấu tính quá đi\~{}!'' Ririmu phồng má trêu anh.

``Ít nhất thì anh cũng nên phụ khiêng nó lên chứ!'' Saku-nee quay ra trách móc. Nhưng Onii-sama chỉ nhún vai: ``Rồi rồi, khi nào có kho báu thật thì gọi anh tới nhặt cũng được\~{}''

``Anh không có phần đâu nha\~{}'' Yuika-nee ngân nga, nhìn anh với ánh mắt trêu chọc. Anh lại chỉ lắc đầu: ``Anh còn chả cần nữa là. Mấy đứa lấy hết đi.''

Tôi khẽ thở dài. Mọi người ai cũng phấn khởi, riêng tôi chỉ thấy\dots\ lo lắng.

Khi họ đã đào gần năm trăm mét - thực sự là sâu đến khó tin, thì\dots

*ỤC ỤC*

``Tiếng\dots\ gì thế?'' Saku-nee dừng xẻng, nghiêng đầu.

``Tôi nghĩ là ma đấy\~{}'' Yuika-nee đáp tỉnh bơ, tay vẫn đào như chưa nghe thấy gì.

``\textbf{Chắc chắn là kho báu đang gọi tụi mình rồi!}'' Ririmu-chan hét lên, mắt sáng lấp lánh.

*ỤC ỤC*

Tôi liếc sang Onii-sama. Vẻ mặt anh bỗng thay đổi, nghiêm túc hơn. ``Tiếng này\dots\ giống tiếng nước thì phải?''

Anh vừa nói xong, chiếc điện thoại của anh lập tức phát ra tiếng còi báo động inh ỏi.

``C-Cái gì vậy!?'' tôi giật bắn người.

Anh nhìn tôi, giọng căng thẳng: ``\eng{Game, ông có chắc-}''

Chợt khựng lại. À phải, không có Game-san ở đây. Giờ tôi là trợ lý.

Tôi nhanh chóng quét lại dữ liệu: ``Có! Có dấu hiệu của một mạch nước ngầm ngay dưới khu vực này!''

``Ể!? Nghiêm túc chứ!?''

Onii-sama không đợi tôi trả lời, liền gào lên: ``\textbf{Này mây đứa! Đừng đào nữa! Nguy hiểm đấy!!}''

*ỤC ỤC*

Nhưng ba chị kia vẫn thản nhiên đào như chưa hề nghe gì. Đặc biệt là Ririmu-chan đúng kiểu say mê công việc đến mức mất kiểm soát.

*ỤC ỤC* *RẮC!*

``T-Tiếng gì nữa vậy!?'' chị Gấu trúc khựng lại, mặt căng như dây đàn.

``Chắc là\dots\ cái đó đấy!'' cô bé nhỏ tuổi mừng rỡ. ``Chắc chắn ở dưới đó là kho báu! Ta sắp đào tới rồi!''

Tôi mở to mắt. Rồi trong một lần chớp nhoáng\dots

\textbf{*BÙM!* *PHỤT!*}

\gif{nj-2l}{1}{225}

Một cột nước vọt thẳng lên trời như mạch phun tự nhiên, mạnh tới mức bốc cả ba người họ lên không trung.

``\textbf{WAAAAH\~{} Tôi đang bay\~{}\dots}'' tiếng thét của Yuika-nee vang vọng cả công viên.

``\textbf{Cái quái gì thế này\dots!}'' Saku-nee kêu lên trước khi mất hút giữa trời mây.

``\textbf{Em chỉ muốn tìm kho báu thôi\dots!}'' đó là lời cuối cùng của Ririmu trước khi cô bé cũng bị hút lên.

Ba bóng người bé xíu cứ thế bị hất tung vào không trung, xoáy tít giữa trời chiều, rồi dần dần nhỏ lại như mấy chấm sáng, hòa lẫn vào bầu trời mùa xuân lộng gió.

Tôi đứng đờ người. ``\vie{Cái gì vừa xảy ra vậy\dots}''

Onii-sama ngước nhìn lên trời, giọng mệt mỏi: ``\vie{\dots Anh đã nói rồi mà. Mấy đứa này cũng óc nho không khác gì bên \hll\ cả.}''

Tôi nuốt nước bọt. Trong lòng chỉ còn một suy nghĩ:
``\textit{Có lẽ mình nên gọi đội cứu hộ trước khi họ rơi trúng nóc nhà ai đó\dots}''

%==========================================TRUYỆN PHỤ=========================================
\omk

\pov{-}

Ngay khoảnh khắc cột nước ngầm bắn lên trời, Kuon không còn giữ được vẻ bình tĩnh thường thấy. Cậu quay sang Suzuki, mặt đanh lại.

``\vie{Đi thôi. \emph{Ngay bây giờ.}}''

``\vie{Không chờ họ rơi xuống ạ?}'' Suzuki nhíu mày, mắt ngước lên trời tỏ vẻ lo lắng.

``\vie{Không. Tốt nhất là chúng ta không liên quan gì đến cái bãi lầy pháp lý này.}''

Nói rồi, cả hai cắm đầu chạy khỏi hiện trường với tốc độ đáng kinh ngạc. Cột nước phía sau vẫn tiếp tục phun như một vòi rồng khổng lồ không kiểm soát, bắn tung tóe nước ra mọi hướng. Đất đá quanh gốc cây bị bào mòn, cây anh đào cổ thụ rung bần bật rồi nghiêng hẳn về một bên. Những cánh hoa bị cuốn theo dòng nước như thể đang khóc than cho số phận đám người vừa ``bay màu''.

Người dân sống gần khu vực sông lập tức hoảng hốt. Có người thì giơ điện thoại quay lại cảnh tượng kỳ lạ này. Có người thì chạy tán loạn vì nghĩ rằng  nơi này sắp sụp đổ. Một số cụ già thì thì thầm về điềm xấu hay ``hồn ma trú ngụ dưới gốc cây anh đào'', góp phần biến cảnh tượng kỳ quái này trở thành một\dots\ truyền thuyết đô thị mới nổi chỉ trong vài giờ tiếp theo.

Trong lúc đó, Kuon và Suzuki đã nấp sau một cửa hàng bách hóa ven đường.

``\vie{Ít nhất thì cũng không có cảnh sát nào tới kịp\dots}'' cậu Tổng thở dốc, tóc ướt nhẹp.

Cô em út thở dài, rũ nước ra khỏi tay áo: ``\vie{Chắc chắn là có báo đài đấy. Em nghĩ mấy cái máy quay bay xung quanh không phải đồ chơi đâu.}''

``Kệ! Không liên quan đến mình là được!'' cậu lẩm bẩm như một câu thần chú, rồi thở dài thườn thượt.

Ririmu, với khả năng chống chịu vô lý bẩm sinh, rơi xuống một chiếc lưới căng trên sân thượng công ty một cách an toàn tuyệt đối. Nhưng khi cô bé tỉnh lại, cô không nhớ mình đang ở đâu, chỉ thấy một nhân viên bảo vệ đang chĩa đèn pin vào mặt cô.

Còn Yuika và Saku thì, đáng buồn thay, bay xa đến mức không còn thấy bóng dáng. Đội tìm kiếm chỉ phát hiện được chiếc xẻng của họ nằm cách gốc cây gần một cây số, cùng một tờ ghi chú viết tay nguệch ngoạc: ``Chúng tôi vẫn sống! Hẹn gặp lại sau 10 ngày nữa!''

Cả công ty Nijisanji rối tung cả lên sau vụ đó. Tin tức lan ra khiến các nhân viên kỹ thuật bắt đầu lên kế hoạch kiểm tra lại toàn bộ hệ thống giám sát địa chất xung quanh trụ sở. Họ còn phải gọi thợ hàn để sửa lại hàng rào bị phá nát bởi áp lực nước.

Kuon và Suzuki và thậm chí Kaigou một lần nữa khai rằng không liên quan tới chuyện này.

\ct

Tối hôm đó, trong một góc quán ăn nhỏ của khu phố mua sắm, Kuon và Suzuki ngồi nhấm nháp mì soba, đèn vàng ấm hắt lên gương mặt mệt mỏi.

``Ít nhất thì không ai nhận ra anh là người ngoài hành tinh\dots'' cô bé đùa khẽ.

Kuon cười khan: ``Nếu đây là ngày thứ hai, thì anh không muốn biết những ngày tiếp theo sẽ tệ đến mức nào.''

Cả hai cụng đũa, rồi im lặng nhìn ra cửa sổ. Ngoài kia, đèn phố nhấp nháy, còn cột nước thì vẫn đang phun nhẹ. Hàng người vẫn đứng xung quanh nhìn hiện tượng kỳ quặc kia, vẫn bu đầy ở quanh nó, tạo ra khung cảnh trái ngược của hai thế giới.

\ct

Trời đã ngả về tối khi hai anh em họ Hikawa trở về căn nhà thuê tạm, nơi mà họ, cùng Kaigou và thực thể Game, đang sống trong mười hai ngày lưu đày tạm thời ở thế giới của Nijisanji.

Khi bước qua cánh cửa gỗ cũ, họ được chào đón bởi một khung cảnh mà Kuon không ngờ sẽ có trong ngày hôm nay: phòng khách sạch bóng, gối đệm được xếp ngay ngắn, bát đĩa trong bếp đã được rửa và úp gọn gàng, thậm chí mấy lọ gia vị còn được phân loại lại theo màu.

Ở giữa gian phòng, Kaigou đang ngồi gọt táo, vẻ mặt lạnh lùng nhưng có phần thư giãn. Cô chỉ liếc mắt nhìn họ. ``\vie{Về rồi à.}''

Suzuki tháo giày, lấm lét nhìn xung quanh. ``\vie{Nhà\dots\ sạch thế này\dots}''

Kaigou nhìn họ bằng ánh mắt bình thản, như thể vừa quét sạch một đám tàn quân xâm lược.

``\vie{Lúc đi chợ, chị nghe nói có cái cột nước tự nhiên phun lên ở công viên. Cảnh sát còn phải phong tỏa khu vực,}'' cô dừng lại, lột vỏ miếng táo cuối cùng rồi nhìn thẳng vào mắt Kuon. ``\vie{Hai người có biết gì về chuyện đó không?}''

Cậu lập tức quay đi, làm như đang tìm giẻ lau trong ba-lô: ``\vie{Không, tôi không thấy gì cả.}''

``\vie{Chắc là ai đó nghịch dại thôi,}'' Suzuki luống cuống. ``\vie{Chị biết đấy\dots\ người ở đây kỳ lắm.}''

``\vie{Phải, phải,}'' Kuon gật gù, ``\vie{có thể là do hệ thống cũ. Thỉnh thoảng nước cũng tự phun mà, ha ha\dots}''

Kaigou nhìn cả hai bằng ánh mắt sắc như dao thái sashimi, nhưng rốt cuộc chỉ hừ nhẹ một tiếng, cười nhẹ. ``\vie{Ừ. Biết vậy là được rồi. Lần sau nhớ mang muối tiêu khi về. Tôi cần chúng cho món trưa mai.}''

``V-Vâng!/Được rồi\dots'' cả hai đồng thanh.

Cô nàng đứng dậy, cầm táo đã gọt xong đi về phía bếp. Khi vừa khuất sau cánh cửa, Suzuki thì thầm: ``\vie{\hl{C-Có khi Onee-sama biết hết rồi\dots}}''

Kuon chỉ kéo dài tiếng thở dài, rồi ngồi phịch xuống đệm: ``\vie{Kaigou-chan tưởng thế mà tinh thật. Thôi thì, ít nhất cô ấy vẫn để cho anh em mình sống\dots}''

Với anh em nhà Hikawa, chuyện như thế là quá đủ cho ngày thứ hai ở đây.

\end{document}